\documentclass[11pt]{article}
\usepackage{times}
\usepackage{fullpage}
\usepackage{url}
\usepackage{hyperref}
\usepackage{fancyhdr}
\usepackage{graphicx}
\usepackage{enumitem}
\usepackage{subcaption}
\usepackage{amsmath, amsfonts, amsthm, fouriernc}
\usepackage{IEEEtrantools}
\usepackage{parskip}

\pagestyle{fancy}
\addtolength{\headheight}{2em}
\addtolength{\headsep}{1.5em}
\lhead{\large{ME 2801}}

\usepackage{sectsty}
\sectionfont{\fontsize{12}{8}\selectfont}

\begin{document}

\begin{center}
  \Large{\bf{Assignment: Stability Margins and Gain Design}}
\end{center}

The textbook exercises designated with the word ``Nise'' are from the
``Problems'' section at the end of Chapters~10 and~11 in \emph{Control System
Engineering} by Norman Nise, 7th edition.

% ---------------------------------------------------------------
\section*{Part 1: Reading Stability Margins from a Bode Plot}

\begin{enumerate}

\item For each open-loop transfer function, use MATLAB to plot the Bode diagram
  and read off the gain margin (GM) and phase margin (PM).  State whether
  the closed-loop system is stable and explain how you determined this from
  the margins.

  \begin{enumerate}
    \item $G(s) = \dfrac{100}{(s+1)(s+5)(s+10)}$
    \item $G(s) = \dfrac{10}{s(s+2)(s+8)}$
    \item $G(s) = \dfrac{K}{s(s+1)(s+4)}$ with $K = 1$
  \end{enumerate}

\item For part~(c) above, use MATLAB \texttt{margin()} to confirm your
  margin values.  Then determine the maximum gain $K_{\max}$ that keeps the
  closed loop stable.

\end{enumerate}

% ---------------------------------------------------------------
\section*{Part 2: Relating Frequency Response to Transient Response}

\begin{enumerate}[resume]

\item A unity-feedback system has open-loop transfer function
  \[
    G(s) = \frac{K}{s(s+4)}
  \]
  \begin{enumerate}
    \item For $K = 4$, compute the phase margin using MATLAB.
    \item Use the approximation $\zeta \approx \mathrm{PM}/100$ (for
      $\mathrm{PM} < 70°$) to estimate the damping ratio of the dominant
      closed-loop poles.
    \item Use MATLAB \texttt{step()} on the \emph{closed-loop} system to find
      the actual \%OS, and compare it to the prediction
      $\%OS \approx e^{-\pi\zeta/\sqrt{1-\zeta^2}} \times 100$.
  \end{enumerate}

\item \textbf{Nise 10.7} (selected parts): Read gain and phase margins from a
  given Bode plot and determine closed-loop stability.

\end{enumerate}

% ---------------------------------------------------------------
\section*{Part 3: Gain Adjustment Design}

\begin{enumerate}[resume]

\item A unity-feedback position-control system has open-loop transfer function
  \[
    G(s) = \frac{K}{s(s+2)(s+5)}
  \]
  Design a gain $K$ such that the closed-loop phase margin is at least $40°$.

  \begin{enumerate}
    \item Plot the Bode diagram for $K = 1$ using MATLAB.
    \item Identify the frequency $\omega_{PM}$ at which the phase equals
      $-180° + 40° = -140°$.
    \item Compute the required gain adjustment so the magnitude is $0$\,dB at
      $\omega_{PM}$.
    \item Verify by re-plotting with the designed $K$ and confirming the margin
      with \texttt{margin()}.
  \end{enumerate}

\item \textbf{Nise 11.2} (selected parts): Given a desired percent overshoot,
  find the gain $K$ using the phase-margin/damping-ratio relationship.

\end{enumerate}

% ---------------------------------------------------------------
\section*{Part 4: USV Application}

\begin{enumerate}[resume]

\item For the USV surge plant with proportional control:
  \[
    L(s) = K_p \cdot G_{\mathrm{surge}}(s) = \frac{5.1\,K_p}{1.3\,s + 1}
  \]
  \begin{enumerate}
    \item Show that this first-order loop has infinite gain margin and
      compute the phase margin as a function of $K_p$.
      (Hint: a first-order system's phase never reaches $-180°$.)
    \item What does this imply about the stability of proportional control of
      a first-order plant?
    \item Verify with MATLAB \texttt{margin()} for $K_p = 2$.
  \end{enumerate}

\end{enumerate}

\end{document}
